\chapter{Estado del arte}
\label{cap:estado}

Este capítulo revisa los principales enfoques de la literatura científica y técnica en torno a agentes de lenguaje aplicados al desarrollo de software, marcos multi-agente conversacionales, orquestación mediante grafos de tareas y políticas de optimización de costos. Finalmente, se presenta una matriz comparativa que posiciona las contribuciones de \sistema\ frente a las limitaciones de los sistemas existentes.

\section{Agentes autónomos individuales en ingeniería de software}

El surgimiento de conjuntos de evaluación rigurosos construidos sobre repositorios reales de código abierto, notablemente SWE-bench \cite{jimenez2024swebench}, impulsó una nueva generación de agentes de software. SWE-bench recopila más de dos mil incidencias históricas resueltas en proyectos populares de Python en GitHub, exigiendo a los modelos generar parches que superen las pruebas de regresión del proyecto sin acceso a las soluciones humanas.

En este contexto, SWE-agent \cite{yang2024sweagent} demostró que la efectividad de un modelo de lenguaje depende fuertemente de la \emph{Interfaz Agente-Computadora} (\emph{Agent-Computer Interface}, ACI). Al proporcionar al modelo comandos adaptados para navegar directorios, realizar búsquedas de texto paginadas y editar bloques acotados mediante coincidencia exacta de líneas, SWE-agent logró mejoras dramáticas de desempeño frente a interfaces basadas en shells estándar. Herramientas contemporáneas orientadas al flujo de trabajo del desarrollador en terminal —tales como Aider, OpenHands y las CLI especializadas Claude Code o Codex CLI \cite{chen2021codex}— han profundizado este enfoque, incorporando mapas de repositorio basados en árboles sintácticos (Tree-sitter) y bucles interactivos de edición y prueba.

A pesar de su notable efectividad en la resolución de incidencias autocontenidas o refactorizaciones locales, estas herramientas operan fundamentalmente bajo el paradigma de \textbf{un agente único por problema}: el modelo ejecuta un bucle lineal de inferencia hasta agotar su presupuesto de tokens o turnos. Como se evidenció en la motivación, este enfoque monolítico no escala cuando el requerimiento exige modificaciones arquitectónicas simultáneas en múltiples subsistemas o capas (por ejemplo, esquemas de persistencia, reglas de dominio, controladores de API e interfaces de usuario), ya que la masa de código excede la ventana de contexto efectiva y desencadena alucinaciones, pérdida de dependencias y desbordes de alcance. \sistema\ no compite como un modelo generativo individual frente a estas herramientas, sino que actúa como un \textbf{plano supervisor de control, partición y verificación} que las coordina concurrentemente.

\section{Sociedades multi-agente conversacionales}

Para superar las limitaciones del agente individual, diversos trabajos han propuesto marcos basados en \emph{sociedades multi-agente} con división de roles. ChatDev \cite{qian2024chatdev} y MetaGPT \cite{hong2024metagpt} estructuran la colaboración emulando el modelo en cascada de una empresa de software: definen agentes con perfiles especializados (Product Manager, Arquitecto, Ingeniero de Software, Tester) que se comunican entre sí mediante lenguaje natural. MetaGPT introduce \emph{Procedimientos Operativos Estandarizados} (\emph{Standard Operating Procedures}, SOPs), exigiendo que los agentes intercambien diagramas de clases y documentos de diseño estructurados para reducir la ambigüedad.

Sin embargo, estos marcos conversacionales presentan tres limitaciones severas cuando se aplican a repositorios de software reales preexistentes:

\begin{enumerate}
  \item \textbf{Generación desde cero vs. evolución de código existente:} Estos sistemas fueron diseñados principalmente para generar aplicaciones de juguete desde cero (*greenfield*). Carecen de mecanismos para indexar repositorios complejos con millones de líneas de código, historiales de Git o configuraciones de compilación preexistentes.
  \item \textbf{Deriva conversacional y desalineación de interfaces:} Al basar la coordinación en el intercambio libre de mensajes de texto entre modelos, el código producido por el agente programador frecuentemente contradice las signaturas asumidas por el agente tester o el arquitecto, generando bucles infinitos de discusión sin convergencia matemática.
  \item \textbf{Ausencia de control de versiones y contención física:} No implementan aislamiento a nivel de sistema operativo ni control de versiones por intento. Las escrituras se realizan de forma directa y destructiva en el sistema de archivos, sin capacidad de revertir intentos fallidos ni auditar los cambios contra un árbol de Git limpio.
\end{enumerate}

\section{Orquestación de tareas mediante grafos}

La descomposición de flujos de razonamiento en grafos dirigidos acíclicos ha sido explorada por sistemas como LLMCompiler \cite{kim2024llmcompiler} y DSPy. LLMCompiler analiza las dependencias de datos entre llamadas a funciones y construye un DAG de invocaciones, despachando en paralelo aquellas operaciones cuyas entradas no dependen de los resultados de otras herramientas.

Si bien este principio de paralelización topológica es conceptualmente valioso, en el dominio de la ingeniería de software las dependencias no se limitan a valores escalares o salidas JSON retornadas por una API:
\begin{itemize}
  \item Las dependencias en software son \textbf{parches de código sobre el sistema de archivos}, los cuales pueden generar conflictos de fusión textual o semántica si dos agentes editan archivos concurrentes.
  \item Requieren \textbf{contratos de interfaces públicas congeladas (\emph{seams})} para habilitar trabajo paralelo seguro antes de que la implementación exista.
  \item Exigen \textbf{obligaciones jerárquicas de validación} para demostrar que la composición de dos módulos preserva el comportamiento global.
\end{itemize}

Por ende, un DAG de herramientas homogéneo no puede resolver la coordinación de agentes de código sin un sistema formal de contratos y relaciones tipadas.

\section{Políticas de asignación de recursos y enrutamiento}

La optimización de costos y recursos en el uso de modelos de lenguaje ha sido abordada mediante técnicas de \emph{enrutamiento adaptativo de modelos}. FrugalGPT \cite{chen2023frugalgpt} y RouteLLM \cite{ong2024routellm} entrenan clasificadores que analizan la complejidad de una consulta y la redirigen al modelo más económico capaz de resolverla (por ejemplo, derivando consultas simples a modelos ligeros como GPT-4o-mini y reservando modelos de mayor capacidad como Claude 3.5 Sonnet para tareas complejas).

Estos enfoques resuelven la pregunta de \textbf{<<qué modelo asignar a una consulta dada>>}. Por el contrario, la política de granularidad desarrollada en esta tesis actúa en una etapa previa y complementaria, resolviendo la pregunta de \textbf{<<cuánto trabajo constituye una unidad atómica ejecutable>>} antes de la asignación. Ambas estrategias son ortogonales y pueden combinarse armónicamente.

\section{Análisis de brechas y posicionamiento de ManyHands}

La tabla~\ref{tab:estado-del-arte} resume la comparación sistemática entre los enfoques representativos del estado del arte y la arquitectura de \sistema.

\begin{table}[htbp]
  \centering
  \caption{Comparación sistemática de ManyHands frente al estado del arte.}
  \label{tab:estado-del-arte}
  \small
  \begin{tabularx}{\textwidth}{@{}l X X X X@{}}
    \toprule
    \textbf{Dimensión} & \textbf{SWE-agent} \cite{yang2024sweagent} & \textbf{MetaGPT} \cite{hong2024metagpt} & \textbf{LLMCompiler} \cite{kim2024llmcompiler} & \textbf{ManyHands} (Esta tesis) \\
    \midrule
    \textbf{Unidad de trabajo} & Incidencia aislada (1 agente) & Proyecto completo (roles conversacionales) & DAG de llamadas a herramientas & \textbf{Run de producto con grafo híbrido} \\
    \textbf{Descomposición} & No (monolítico) & Fija por roles (PM, Dev, QA) & Dinámica por flujo de datos & \textbf{Adaptativa categórica (Política 4.0)} \\
    \textbf{Fronteras de interfaz} & N/A & Prompts en lenguaje natural & Parámetros de función & \textbf{Contratos formales versionados (\emph{seams})} \\
    \textbf{Aislamiento físico} & Entorno local único & Directorio compartido & N/A (en memoria) & \textbf{Pool de Git worktrees con scope guard} \\
    \textbf{Resolución de conflictos} & N/A & Discusión multi-turno & N/A & \textbf{Integración bottom-up con repair} \\
    \textbf{Verificación} & Test final de la incidencia & Generación de tests por LLM & N/A & \textbf{Matriz de Evidencias sobre commit exacto} \\
    \textbf{Persistencia} & Logs de texto & Logs de texto & N/A & \textbf{Journal append-only (Event Sourcing)} \\
    \textbf{Entrega} & Parche \texttt{.patch} & Archivos generados & Resultado en memoria & \textbf{Manifiesto y recibo transaccional} \\
    \bottomrule
  \end{tabularx}
\end{table}

A partir de este análisis de brechas, las contribuciones distintivas de \sistema\ se sintetizan en tres pilares de ingeniería:

\begin{enumerate}
  \item \textbf{Gobernanza mediante contratos y relaciones tipadas:} Sustitución de la comunicación libre entre agentes por un plano formal de contratos inmutables (\texttt{TaskContractBundle}, \texttt{SeamContract}, \texttt{ArtifactContract}, \texttt{ValidationObligation}) que garantiza la compatibilidad de interfaces antes de la ejecución.
  \item \textbf{Política de granularidad categórica explicable:} Superación de los modelos opacos de scoring y las reglas estáticas mediante una política que evalúa capacidad física, paralelismo real y aislamiento de pruebas, respaldada por un banco de regresión offline congelado.
  \item \textbf{Disciplina de evidencia y entrega transaccional:} Validación estricta sobre el SHA del commit exacto en worktrees aislados, con auditoría de controles negativos e integridad de tests, publicando únicamente resultados con recibos inmutables de entrega.
\end{enumerate}
